\documentclass[10pt,a4paper]{amsart}
\usepackage[margin=19mm]{geometry}
\usepackage{amsmath,amssymb,mathtools}
\usepackage[scr=rsfs,cal=euler]{mathalfa}
\usepackage{xfrac}
\usepackage[unicode,hidelinks,bookmarks=false]{hyperref}
\newcommand{\ie}{i.e.\ }
\newcommand{\cf}{cf.\ }
\newcommand{\resp}{respectively\ }
\newcommand{\Subsection}{\subsection}
\providecommand{\nobreakdash}{\nobreak\hbox{-}}
\setlength{\emergencystretch}{2em}
\multlinegap0pt
\usepackage{bm}
\usepackage{bbm}
\usepackage{dsfont}
\usepackage{xspace}
\usepackage{cancel}
\usepackage{mathtools}
\usepackage{amsthm}
\makeatletter
\@ifundefined{lemma}{\newtheorem{lemma}{Lemma}}{}
\makeatother
\title[How many points does an affine algebraic set have in residue]{How many points does an affine algebraic set have in residue classes modulo n ?}
\author{Mihai Prunescu}
\date{Source: arXiv:2608.22049v4 (CC BY 4.0)}
\begin{document}
\maketitle

\noindent\emph{Source and licence.} How many points does an affine algebraic set have in residue classes modulo n ?. Version arXiv:2608.22049v4 (CC BY 4.0), distributed under the Creative Commons Attribution 4.0 International licence, as recorded in the arXiv licence metadata for this exact version and shown on the abstract page https://arxiv.org/abs/2608.22049v4. Full rights notice: licenses/arxiv-2608.22049-CC-BY-4.0.txt.\par\smallskip

\noindent\textit{Editorial scope.} Counted unit: the Lemma whose source label is \texttt{lem:sufficientt} in the article (source0.tex lines 252--255, source label \texttt{lem:sufficientt}), delivered together with the article's own complete original proof of it (source0.tex lines 257--291, one \texttt{proof} environment of 35 lines). No mathematics is written, repaired or re-derived: statement and proof are the article's own text line for line. Required context retained uncounted: none. Every definition, display and label of those lines is retained unchanged. Omitted from this package and not redistributed: 1--251, 256--256, 292--392. Nothing inside a retained excerpt is omitted or altered: every retained body is delivered exactly as the article's lines are written, so no omission receipt of any kind accompanies this package. Citations: the retained excerpts carry no citation command at all, so no bibliography entry of the article is transcribed and no reference is redistributed. Reference adaptation: none was needed and none was made; every reference inside the delivered unit resolves inside it, so no unresolved reference or citation is produced. Printed slips kept literal: no printed word, symbol, hypothesis or number of the counted statement or of its proof was altered. Layout and class: the article's own class is \texttt{article} with packages including a4wide, amsmath, amssymb, amsthm, inputenc, babel, natbib, url, color, hyperref, cleveref, tabularx, breqn, appendix; the wrapper is the lane's pinned \texttt{amsart} editorial wrapper at 10pt on A4, which restates the theorem-like environments the retained text needs and adds the article's own macro definitions that the retained text needs (none), with extra wrapper packages (bm, bbm, dsfont, xspace, cancel, mathtools). The delivered numbering is the unit's own. Assets: no retained span contains a figure environment, a graphics-inclusion command, a PDF-page inclusion, a table or a TikZ picture; no image file is copied into this package, the package declares no asset dependency and no image file is redistributed with it. Licence: version arXiv:2608.22049v4 of this article is distributed under the Creative Commons Attribution 4.0 International licence, as recorded in the arXiv licence metadata for this exact version and shown on the abstract page https://arxiv.org/abs/2608.22049v4. A full rights notice is delivered at licenses/arxiv-2608.22049-CC-BY-4.0.txt. Characters: every retained excerpt is pure ASCII, so the delivered engine, XeLaTeX with its default Latin Modern text font, reports no missing character.
\begin{lemma}\label{lem:sufficientt}
   There is an arithmetic term $t(n)$ such that given the system $S$, for all $\vec \lambda \in \mathbb Z^s$ and for all $n \geq 1$, 
   $$U_n(\vec \lambda ) \subset [0, t(n)]^{2m + k}.$$
\end{lemma} 

\begin{proof} In the equation $G(\vec \lambda, n, \vec x, \vec y, \vec z) = 0$, it is clear that the $x$-components and the $z$-components of every solution satisfy $x_j < n$ and $z_j < n$. We look now at the components $y_i$. We know that
$$y_i = \frac{|G_i(\vec \lambda, n, \vec x)|}{n}.$$
We perform the following actions on $G_i(\vec \lambda, n, \vec x)$:
\begin{enumerate}
    \item All subtraction signs are replaced with addition signs.
    \item All coefficients $(c \bmod n)$ are replaced by $n$, as we know that $0 \leq c \bmod n < n$.
    \item All occurrences of a parameter in the form $(\lambda_{e} \bmod n) $ are replaced by $n$, as we know that $0 \leq \lambda_e \bmod n < n$.
    \item All variables $x_j$ are replaced by $n$, as we know that as components of a natural solution, one has $x_j < n$. 
\end{enumerate} 
The result is a polynomial $H_i(n)$ with the property that $n \,|\,H_i(n)$ and
$$y_i \le \frac{H_i(n)}{n},$$
for every component $y_i$ of a solution. Now one chooses $t(n)$ equal to the sum of all these bounds:
$$t(n) = \frac{H_1(n)}{n}+ \dots + \frac{H_k(n)}{n} + n.$$
Of course $t(n)$ can be taken $1 + \max(H_1(n), \dots, H_k(n), n)$ and the $\max$ operation can be also expressed by an arithmetic term, but this expression is more complicated. 
\qed 

\begin{lemma}
    Given the system $S$, there is a constant $K \in \mathbb N$ such that for all parameters $\vec \lambda_1, \dots, \vec \lambda_k \in \mathbb Z$ and for all $n \geq 1$, 
    $$0 \leq G(\vec \lambda, n, \vec x, \vec y, \vec z) < 2^{n+K},$$
    for all $(\vec x, \vec y, \vec z) \in [0, t(n)]^{2m + k}$. 
\end{lemma} 

{\bf Proof}: For the function $G(\vec \lambda, n, \vec x, \vec y, \vec z)$ we apply almost the same steps as in the Lemma \ref{lem:sufficientt}, as follows: 
\begin{enumerate}
    \item All subtraction signs are replaced with addition signs.
    \item All coefficients $(c \bmod n)$ are replaced by $n$.
    \item All occurrences of a parameter $(\lambda_{e} \bmod n) $ are replaced by $n$.
    \item All variables $x_j$, $y_i$ and $z_j$ are replaced by $t(n)$. 
\end{enumerate} 
The result is a polynomial $H(n)$ such that for all parameters $\vec \lambda \in \mathbb Z^s$, for all $n \geq 1$, and for all $(\vec x, \vec y, \vec z) \in [0, t(n)]^{2m+k}$, 
$$0 \leq G(\vec \lambda,n, \vec x, \vec y, \vec z) \leq H(n). $$
But as $H(n)$ is a polynomial, there is always a constant $K$ such that for all $n \in \mathbb N$, 
$$H(n) < 2^{n+K}.$$
Indeed, $\lim_{n \rightarrow \infty} 2^n / |H(n)| = + \infty$, so there is some $N_0$ such that the inequality is surely true for $n \geq N_0$. For sufficiently large $K$, $2^K \times 2^n > H(n)$ for all $n \in \mathbb N$. 
\end{proof}

\end{document}
